\beginnumbering
\eledsection*{Adhyāya 1}
\pstart \ledsidenote{\passref{1.1.0}}\edtext{\skt{नमः}}{\Afootnote{N K; \skt{ओं नमो}\ H}}\skt{ }\edtext{\skt{कमलहस्ताय}}{\Afootnote{K; \skt{नागार्जुनपादाय}\ N; \skt{धन्वन्तरये}\ H}}\skt{ ॥}\pend

\pstart \ledsidenote{\passref{1.1.1}}\skt{अथातो }\edtext{\skt{वेदोत्पत्तिमध्यायं}}{\lemma{\skt{अध्यायं}}\Afootnote{A K\textsuperscript{pc}; \skt{आध्यायं}\ K\textsuperscript{ac} N; \skt{नामाध्यायं}\ H}}\skt{ व्याख्यास्यामः ।}\pend

\pstart \ledsidenote{\passref{1.1.3}}\skt{अथ खलु }\edtext{\skt{भगवन्तममरवरमृषिगणपरिवृत्तमाश्रमस्थं}}{\lemma{\skt{अमरवरम्}}\Afootnote{A K H; \skt{अमरवर}\ N}}\skt{ काशिराजं }\edtext{\edtext{\skt{दिवोदासमौपधेनववैतरणौरभ्रपुष्कलावतकरवीरगोपुररक्षितभोजसुश्रुतप्रभृतय}}{\lemma{\skt{दिवोदासम्}}\Afootnote{A K\textsuperscript{pc} N H; \skt{दिवोदासम्म्}\ K\textsuperscript{ac}}}}{\lemma{\skt{औपधेनववैतरणौरभ्रपुष्कलावतकरवीरगोपुररक्षितभोजसुश्रुतप्रभृतय}}\Afootnote{\skt{धन्वन्तरिमौपधेनववैतरणौरभ्रपौष्कलावतकरवीर्य(र)गोपुररक्षितसुश्रुतप्रभृतय}\ A; \skt{औपधेनवः वैतरणोरभ्रपुष्कलावतकरवीरगोपुररक्षितभोजसुश्रुतप्रभृतय}\ K\textsuperscript{ac}; \skt{धन्वन्तरिमौपधेनवः वैतरणोरभ्रपुष्कलावतकरवीरगोपुररक्षितभोजसुश्रुतप्रभृतय}\ K\textsuperscript{pc}; \skt{औपधेनववैस्तरणौरभ्रपौष्कलावतकरवीरगोपुररक्षितभोजसुश्रुतप्रभृतव्य}\ N; \skt{धन्वन्तरिमौपधेनववैरणोरभ्रपुष्कलावतकरवीरगोपुररक्षितभोजसुश्रुतप्रभृतय}\ H}}\footnoteA{The scribe of the earliest MS, K, added the name \emph{dhanvantari} in the margin; the scribe of N did not have the name; the scribe of H included it in the text.}\skt{ }\edtext{\skt{ऊचुः}}{\Afootnote{A N H; \skt{ऊचु}\ K}}\skt{ ।}\pend

\pstart \ledsidenote{\passref{1.1.4}}\edtext{\edtext{\edtext{\edtext{\skt{भगवञ्शारीरमानसागन्तुभिर्व्याधिभिर्विविधवेदनाभिघातोपद्रुतान्सनाथाननाथवद्विचेष्टमानान्विक्रोशतश्च}}{\lemma{\skt{भगवञ्शारीरमानसागन्तुभिर्}}\Afootnote{\skt{भगवन्शारीरमानसागन्तुभिर्}\ A; \skt{भगवरीरमानसागन्तु भि…}\ K\textsuperscript{ac}; \skt{भगवच्छारीरमानसागन्तु भि…}\ K\textsuperscript{pc}; \skt{भगवञ्छरीरमानसागन्तुभिर्}\ N; \skt{भगवच्छारीरमानसागन्तुभिर्}\ H}}}{\lemma{\skt{व्याधिभिर्}}\Afootnote{A K H; \skt{व्याधिभि}\ N}}}{\lemma{\skt{विविधवेदनाभिघातोपद्रुतान्सनाथाननाथवद्विचेष्टमानान्विक्रोशतश्}}\Afootnote{\skt{विविधवेदनाभिघातोपद्रुतान्सनाथानप्यनाथवद्विचेष्टमानान्विक्रोशतश्}\ A; \skt{वि…धवेदनाभिघातोपद्रुतां सनाथाननाथवद्विचेष्टमानान्विक्रोशत}\ K; \skt{विविधवेदनाभिघातोपद्रुतानासनाथाननाथवद्विचेष्टमानान्विक्रोशतश्}\ N; \skt{व्विविधवेदनाभिघातोपद्रुतां सनाथाननाथवद्विचेष्टमानान्विक्रोशतश्}\ H}}}{\lemma{\skt{च}}\Afootnote{A N H; \skt{…श्च}\ K}}\skt{ }\edtext{\skt{मानवानभिसमीक्ष्य}}{\Afootnote{A K\textsuperscript{ac} N; \skt{मानवानमभिसमीक्ष्य}\ K\textsuperscript{pc}; \skt{मानवानामभिसमीक्ष्य}\ H}}\skt{ मनसि नः }\edtext{\skt{पीडाभवत्}}{\Afootnote{K N H; \skt{पीडा भवति}\ A}}\skt{ ।}\pend

\pstart \ledsidenote{\passref{1.1.4b}}\skt{तेषां सुखैषिणां }\edtext{\skt{रोगोपशमार्थमात्मनश्च}}{\lemma{\skt{रोगोपशमार्थम्}}\Afootnote{A K H; \skt{रोगोपगमार्थम्}\ N}}\skt{ }\edtext{\edtext{\skt{प्राणयात्रार्थमायुर्वेदं}}{\lemma{\skt{प्राणयात्रार्थम्}}\Afootnote{A K H; \skt{प्राणयातार्थम्}\ N}}}{\lemma{\skt{आयुर्वेदं}}\Afootnote{N H; \skt{प्रजाहितहेतोरायुर्वेदं}\ A; \skt{आयु}\ K}}\skt{ }\edtext{\skt{इच्छाम}}{\Afootnote{N H; \skt{श्रोतुमिच्छाम}\ A; \skt{श्वेदमिच्छाम}\ K}}\skt{ }\edtext{\skt{उपदिश्यमानम्}}{\Afootnote{K N H; \skt{इहोपदिश्यमानम्}\ A}}\skt{ । }\edtext{\edtext{\skt{अत्रायत्तमैहिकमामुष्मिकञ्च}}{\lemma{\skt{अत्रायत्तम्}}\Afootnote{A K N; \skt{अत्त्राय उत्तम्}\ H}}}{\lemma{\skt{ऐहिकम्}}\Afootnote{A K N; \skt{औहिकम्}\ H}}\skt{ }\edtext{\edtext{\skt{श्रेयस्तद्भगवन्तमुपसन्नाः}}{\lemma{\skt{तद्भगवन्तम्}}\Afootnote{K N H; \skt{तद्भगवन्तम्}\ A}}\skt{ स्मः}}{\lemma{\skt{उपसन्नाः स्मः}}\Afootnote{\skt{उपपन्नाः स्मः}\ A; \skt{उपसन्ना स्म}\ K N; \skt{उपसन्ना स्मः}\ H}}\skt{ }\edtext{\skt{शिष्यत्वेनेति}}{\Afootnote{A N H; \skt{शिष्यत्वे}\ K}}\skt{ ।}\pend

\pstart \ledsidenote{\passref{1.1.5}}\skt{तानुवाच भगवान्स्वागतम्वः }\edtext{\skt{सर्व}}{\Afootnote{A K H; \skt{सव}\ N}}\skt{ एवामीमांस्या }\edtext{\skt{अध्याप्याश्च}}{\Afootnote{A K\textsuperscript{pc} N H; \skt{अ प्याश्च}\ K\textsuperscript{ac}}}\skt{ }\edtext{\skt{भवन्तो}}{\Afootnote{A K N; \skt{भगवन्तो}\ H}}\skt{ वत्साः ।}\pend

\pstart \ledsidenote{\passref{1.1.6}}\skt{इह }\edtext{\skt{खल्वायुर्वेदो}}{\Afootnote{K N H; \skt{खल्वायुर्वेदं}\ A}}\skt{ }\edtext{\skt{नाम}}{\Afootnote{K N H; \skt{नामो}\ A}}\skt{ }\edtext{\edtext{\edtext{\skt{यदुपाङ्गमथर्ववेदस्योक्तमनुत्पाद्यैव}}{\lemma{\skt{यदुपाङ्गम्}}\Afootnote{H; \skt{पाङ्गम्}\ A; \skt{यमुपाङ्गम्}\ K; \skt{यदुपागम्}\ N}}}{\lemma{\skt{अथर्ववेदस्योक्तम्}}\Afootnote{K N H; \skt{अथर्ववेदस्या}\ A}}\skt{ च}}{\lemma{\skt{अनुत्पाद्यैव च}}\Afootnote{K N H; \skt{नुत्पाद्यैव}\ A}}\skt{ प्रजाः }\edtext{\skt{श्लोकशतसहस्रमध्यायसहस्रञ्च}}{\lemma{\skt{श्लोकशतसहस्रम्}}\Afootnote{A K\textsuperscript{pc} N H; \skt{श्लोकसहस्रम्}\ K\textsuperscript{ac}}}\skt{ }\edtext{\edtext{\skt{कृतवान्स्वयम्भूरल्पायुष्काल्पमेधस्त्वञ्चालोक्य}}{\lemma{\skt{अल्पायुष्काल्पमेधस्त्वञ्}}\Afootnote{K N H; \skt{ततोऽल्पायुष्ट्वमल्पमेधस्त्वं}\ A}}}{\lemma{\skt{चालोक्य}}\Afootnote{A K N; \skt{च अलोक्य}\ H}}\skt{ नराणाम्भूयोऽष्टधा }\edtext{\skt{प्रणीतवान्}}{\Afootnote{A N H; \skt{प्रणीतवान्तद्यथा}\ K}}\skt{ ।}\pend

\pstart \ledsidenote{\passref{1.1.7}}\edtext{\skt{तद्यथा}}{\Afootnote{A N H; \emph{om.}\ K}}\skt{ शल्यं शालाक्यं }\edtext{\skt{कायचिकित्सा}}{\Afootnote{A N H; \skt{का…चिकित्सा}\ K}}\skt{ भूतविद्या }\edtext{\skt{कौमारभृत्यमगदतन्त्रं}}{\lemma{\skt{अगदतन्त्रं}}\Afootnote{A K N; \skt{अगदतन्त्र}\ H}}\skt{ रसायनतन्त्रं }\edtext{\edtext{\skt{वाजीकरणतन्त्रमिति}}{\lemma{\skt{वाजीकरणतन्त्रम्}}\Afootnote{A N H; \skt{व्याजीकरणतन्त्रम्}\ K}}}{\lemma{\skt{इति}}\Afootnote{A N H; \skt{इति । …}\ K}}\skt{ ।}\pend

\pstart \ledsidenote{\passref{1.1.8.1}}\edtext{\skt{अथास्य}}{\Afootnote{A N H; \skt{स्य}\ K}}\skt{ }\edtext{\skt{प्रत्येकाङ्गलक्षणसमासः}}{\Afootnote{K N H; \skt{प्रत्यङ्गलक्षणसमासः}\ A}}\skt{ ।}\pend

\pstart \ledsidenote{\passref{1.1.8.1a}}\skt{तत्र }\edtext{\skt{शल्यन्नाम}}{\lemma{\skt{शल्यन्}}\Afootnote{A H; \skt{शंल्यन्}\ K N}}\skt{ }\edtext{\skt{विविधतृणकाष्ठपाषाणपांसुलोहलोष्टास्थिबालनखपूयास्रावदुष्टव्रणान्तर्गर्भशल्योद्धरणार्थं}}{\Afootnote{K; \skt{विविधतृणकाष्ठपाषाणपांशुलोहलोष्टास्थिबालनखपूयास्रावदुष्टव्रणान्तर्गर्भशल्योद्धरणार्थं}\ A; \skt{विविधतृणकाष्ठपाषाणपांसुलोहलोष्टास्थिबालनपूयास्रावदु…णान्…र्गर्भशल्योद्धरणार्थं}\ N; \skt{विविधतृणकाष्ठपाषाणपांशुलोष्टास्थिबालनखपूयाश्रावदुष्टव्रणान्तर्गर्भशल्योद्धरणार्थं}\ H}}\skt{ }\edtext{\skt{यन्त्रशस्त्रक्षाराग्निप्रणिधानविनिश्चयार्थञ्च}}{\Afootnote{K N H; \skt{यन्त्रशस्त्रक्षाराग्निप्रणिधानव्रणविनिश्चयार्थं च (१)}\ A}}\skt{ }\edtext{\edtext{\skt{षष्ट्याभिधानैरिति}}{\lemma{\skt{षष्ट्याभिधानैर्}}\Afootnote{N H; \emph{om.}\ A; \skt{षष्ट्याभिध्…नैर्}\ K}}}{\lemma{\skt{इति}}\Afootnote{K N H; \emph{om.}\ A}}\skt{ ।}\pend

\pstart \ledsidenote{\passref{1.1.8.2}}\edtext{\edtext{\skt{शालाक्यतन्त्रन्नामोर्द्ध्वजत्रुगतानां}}{\lemma{\skt{शालाक्यतन्त्रन्}}\Afootnote{K N H; \skt{शालाक्यं}\ A}}}{\lemma{\skt{नामोर्द्ध्वजत्रुगतानां}}\Afootnote{A K N; \skt{नामोर्ध्वजात्रुगतानां}\ H}}\skt{ }\edtext{\skt{विकाराणां}}{\Afootnote{K H; \emph{om.}\ A; \skt{विकाराणा}\ N}}\skt{ }\edtext{\skt{श्रवणनयनवदनघ्राणादिसंश्रितानां}}{\Afootnote{A H; \skt{श्रवणनयनवदनघ्राणादिसं…स्थितानां}\ K; \skt{श्रवणनयनवदनघ्राणादिसंश्रिताना}\ N}}\skt{ }\edtext{\edtext{\skt{विकाराणामुपशमकरणार्थम्}}{\lemma{\skt{विकाराणाम्}}\Afootnote{K N H; \skt{व्याधीनाम्}\ A}}}{\lemma{\skt{उपशमकरणार्थम्}}\Afootnote{K N H; \skt{उपशमनार्थम् (२)}\ A}}\skt{ ।}\pend

\pstart \ledsidenote{\passref{1.1.8.3}}\edtext{\skt{कायचिकीत्सा}}{\Afootnote{K N H; \skt{कायचिकित्सा}\ A}}\skt{ नाम }\edtext{\skt{सर्वशरीरावस्थितानां}}{\Afootnote{H; \skt{सर्वाङ्गसंश्रितानां}\ A; \skt{सर्वशरीरावस्थितानाञ्च}\ K\textsuperscript{ac} N; \skt{सर्वशरीरावस्थिता}\ K\textsuperscript{pc}}}\skt{ }\edtext{\edtext{\skt{व्याधीनामुपशमकरणार्थं}}{\lemma{\skt{व्याधीनाम्}}\Afootnote{A K\textsuperscript{ac} N H; \skt{धीनाम्}\ K\textsuperscript{pc}}}}{\lemma{\skt{उपशमकरणार्थं}}\Afootnote{K N H; \emph{om.}\ A}}\skt{ }\edtext{\edtext{\skt{ज्वरशोफगुल्मरक्तपित्तोन्मादापस्मारप्रमेहातिसारादीनाञ्च}}{\lemma{\skt{ज्वरशोफगुल्मरक्तपित्तोन्मादापस्मारप्रमेहातिसारादीनाञ्}}\Afootnote{K; \skt{ज्वररक्तपित्तशोषोन्मादापस्मारकुष्ठमेहातिसारादीनाम्}\ A; \skt{ज्वरशोफगुल्मरक्तपित्तोन्मादापस्मारप्रमेहातीसारादीनाञ्}\ N H}}}{\lemma{\skt{च}}\Afootnote{K N H; \skt{उपशमनार्थम् (३)}\ A}}\skt{ ।}\pend

\pstart \ledsidenote{\passref{1.1.8.4}}\skt{भूतविद्या नाम }\edtext{\skt{देवगन्धर्वयक्षराक्षसपितृपिशाचविनायकनागग्रहोपसृष्टचेतसां}}{\Afootnote{\skt{देवासुरगन्धर्वयक्षरक्षःपितृपिशाचनागग्रहाद्युपसृष्टचेतसां}\ A; \skt{देवगन्धर्व…क्षराक्षसपितृपिशाचविनायकनागग्रहोपसृष्टचेतसां}\ K; \skt{देवदीनवगन्धर्वयक्षराक्षसपितृपिशाचविनायकनागग्रहोपसृष्टचेतसां}\ N; \skt{देवदानवगन्धर्व्वयक्षराक्षसपितृपिशाचविनायकनागग्रहोपसृष्टचेतसां}\ H}}\skt{ शान्तिकर्मबलिहरणादिग्रहोपशमनार्थं ।}\pend

\pstart \ledsidenote{\passref{1.1.8.5}}\skt{कौमारभृत्यन्नाम }\edtext{\skt{कुमारभरणधात्रीक्षीरदोषसंशोधनार्थं}}{\Afootnote{A K N; \skt{कुमाराभरणधात्रीक्षीरदोषसंशोधनार्थं}\ H}}\skt{ }\edtext{\skt{दुष्टस्तन्यग्रहसमुत्थितानाञ्च}}{\Afootnote{K N H; \skt{दुष्टस्तन्यग्रहसमुत्थानां च}\ A}}\skt{ }\edtext{\skt{व्याधीनामुपशमकरणार्थम्}}{\lemma{\skt{उपशमकरणार्थम्}}\Afootnote{K H; \skt{उपशमनार्थम् (५)}\ A; \skt{उपशमकरणार्थः}\ N}}\skt{ ।}\pend

\pstart \ledsidenote{\passref{1.1.8.6}}\skt{अगदतन्त्रन्नाम }\edtext{\skt{सर्पकीटदष्टविषव्यञ्जनार्थं}}{\Afootnote{K; \skt{सर्पकीटलूतामूषकादिदष्टविषव्यञ्जनार्थं}\ A; \skt{सर्पकीटलूतादष्टसरीसृपविषव्यञ्जनार्थ}\ N; \skt{सर्प्पकीटलूतादष्टसरीसृपविषव्यञ्जनार्थम्}\ H}}\skt{ }\edtext{\skt{विविधविषवेगोपशमनार्थञ्च}}{\lemma{\skt{विविधविषवेगोपशमनार्थञ्}}\Afootnote{N H; \skt{विविधविषसंयोगोपशमनार्थं}\ A; \skt{विविधविषवेगो… शमनार्थञ्}\ K}}\skt{ ।}\pend

\pstart \ledsidenote{\passref{1.1.8.7}}\skt{रसायनतन्त्रन्नाम }\edtext{\edtext{\skt{वयःस्थापनमायुर्मेधाकरणं}}{\lemma{\skt{वयःस्थापनम्}}\Afootnote{H; \skt{व}\ A; \skt{व…}\ K; \skt{वयस्थापनम्}\ N}}}{\lemma{\skt{आयुर्मेधाकरणं}}\Afootnote{N H; \skt{यःस्थापनमायुर्मेधाबलकरं}\ A; \skt{…}\ K}}\skt{ }\edtext{\skt{व्याध्युपशमकरणार्थञ्च}}{\lemma{\skt{व्याध्युपशमकरणार्थञ्}}\Afootnote{N H; \skt{रोगापहरणसमर्थं}\ A; \skt{…पशमकरणार्थञ्}\ K}}\skt{ ।}\pend

\pstart \ledsidenote{\passref{1.1.8.8}}\edtext{\edtext{\skt{वाजीकरणतन्त्रन्नाम}}{\lemma{\skt{वाजीकरणतन्त्रन्}}\Afootnote{A K\textsuperscript{ac} N H; \skt{व्याजीकरणतन्त्रन्}\ K\textsuperscript{pc}}}}{\lemma{\skt{नाम}}\Afootnote{K N H; \skt{नामा}\ A}}\skt{ }\edtext{\skt{स्वल्पदुष्टक्षीणविशुष्करेतसां}}{\Afootnote{K N H; \skt{ल्पदुष्टक्षीणविशुष्करेतसाम्}\ A}}\skt{ }\edtext{\skt{शुक्राप्यायनप्रसादोपचयजनननिमित्तं}}{\Afootnote{K N; \skt{आप्यायनप्रसादोपचयजनननिमित्तं}\ A; \skt{शुक्राप्यायतनप्रसादोपचयजनननिमित्तं}\ H}}\skt{ प्रहर्षजननार्थञ्च ।}\pend

\pstart \ledsidenote{\passref{1.1.9}}\skt{एवमयमायुर्वेदोऽष्टाङ्ग उपदिश्यते । अत्र }\edtext{\skt{कस्मै किं}}{\Afootnote{A N H; \skt{कस्मै कि}\ K}}\skt{ }\edtext{\edtext{\skt{वर्ण्यतामिति}}{\lemma{\skt{वर्ण्यताम्}}\Afootnote{N H; \skt{उच्यताम्}\ A; \skt{…}\ K}}}{\lemma{\skt{इति}}\Afootnote{A N H; \skt{…}\ K}}\skt{ ।}\pend

\pstart \ledsidenote{\passref{1.1.10}}\edtext{\skt{त }\edtext{\skt{ऊचुरस्माकं}}{\lemma{\skt{अस्माकं}}\Afootnote{A N H; \emph{om.}\ K}}}{\lemma{\skt{त ऊचुर्}}\Afootnote{A N H; \emph{om.}\ K}}\skt{ }\edtext{\skt{सर्वमेव }\edtext{\skt{शल्यज्ञानमलङ्कृत्वोपदिशतु}}{\lemma{\skt{अलङ्कृत्वोपदिशतु}}\Afootnote{N H; \skt{मूलं कृत्वोपदिशतु}\ A; \skt{उपदिशतु}\ K\textsuperscript{ac}; \skt{अलंकृ+ उपदिशतु}\ K\textsuperscript{pc}}}}{\lemma{\skt{सर्वमेव शल्यज्ञानम्}}\Afootnote{H; \skt{सर्वेषामेव शल्यज्ञानं}\ A; \skt{शल्यज्ञानम्}\ K\textsuperscript{ac} K\textsuperscript{pc}; \skt{सर्वेषामेव ननमेव शल्यज्ञानम्}\ N\textsuperscript{ac}; \skt{सर्वेषामेव न शल्यज्ञानम्}\ N\textsuperscript{pc}}}\skt{ भगवानिति ।}\footnoteA{We follow the reading of MSS H and N, pace Harimoto 2013, in spite of the fact that \emph{alaṃkṛ[tvā]} is a marginal addition in K. K is badly damaged at this point, so it is not completely clear what happend to the manuscript, although it does look as if the word was not included in the text and was added by a different hand. But H is usually a faithful witness to the same text as K, and in this case even N concurs. The use of \emph{alaṅkṛtvā} in the sense "make accessible" is slightly unusual and probably explains the later banalization to \emph{mūla}.}\pend

\pstart \ledsidenote{\passref{1.1.11}}\edtext{\skt{स }\edtext{\skt{उवाचैवमस्त्विति}}{\lemma{\skt{अस्त्विति}}\Afootnote{A N H; \skt{…}\ K}}}{\lemma{\skt{स उवाचैवम्}}\Afootnote{A H; \emph{om.}\ K; \skt{स उवाच एवंम्}\ N}}\skt{ ।}\pend

\pstart \ledsidenote{\passref{1.1.12}}\edtext{\skt{त }\edtext{\edtext{\skt{ऊचुर्भूयोऽस्माकं}}{\lemma{\skt{भूयो}}\Afootnote{A N H; \skt{{\symfont ‡}}\ K}}}{\lemma{\skt{ऽस्माकं}}\Afootnote{H; \skt{ऽपि भगवन्तम्}\ A; \skt{…कं}\ K; \skt{ऽस्माक स}\ N}}}{\lemma{\skt{त ऊचुर्}}\Afootnote{A N H; \skt{…}\ K}}\skt{ }\edtext{\edtext{\skt{सर्वेषामेवैकमतीनां}}{\lemma{\skt{सर्वेषाम्}}\Afootnote{K H; \skt{अस्माकम्}\ A; \skt{र्वेषांम्}\ N}}}{\lemma{\skt{एवैकमतीनां}}\Afootnote{K H; \skt{एकमतीनां}\ A; \skt{एवैकमतीनानाम्}\ N}}\skt{ मतमभिसमीक्ष्य सुश्रुतो भगवन्तं }\edtext{\skt{प्रक्ष्यति}}{\Afootnote{A K; \skt{प्रत्य}\ N H}}\skt{ । }\edtext{\skt{अस्योपदिश्यमानं}}{\Afootnote{K; \skt{अस्मै चोपदिश्यमानं}\ A; \skt{क्षस्योपदिश्यमानं}\ N H}}\skt{ वयमप्युपधारयिष्यामः ।}\pend

\pstart \ledsidenote{\passref{1.1.13}}\skt{स उवाचैवमस्त्विति ।}\pend

\pstart \ledsidenote{\passref{1.1.14-15}}\edtext{\skt{इह}}{\Afootnote{K N H; \skt{वत्स सुश्रुत! इह}\ A}}\skt{ }\edtext{\skt{खल्वायुर्वेदप्रयोजनं}}{\Afootnote{A K N; \skt{खल्वायुर्व्वेदे प्रयोजनं}\ H}}\skt{ }\edtext{\skt{व्याध्युपसृष्टस्य}}{\Afootnote{K N H; \skt{व्याध्युपसृष्टानां}\ A}}\skt{ व्याधिपरिमोक्षः }\edtext{\edtext{\edtext{\skt{स्वस्थरक्षणञ्चायुरस्मिन्विदन्त्यनेन}}{\lemma{\skt{स्वस्थरक्षणञ्}}\Afootnote{K N H; \skt{स्वस्थस्य रक्षणं}\ A}}}{\lemma{\skt{चायुर्}}\Afootnote{K N H; \skt{च ॥१४॥ ‘आयुर्}\ A}}}{\lemma{\skt{अस्मिन्विदन्त्य्}}\Afootnote{K\textsuperscript{ac} N H; \skt{अस्मिन्विद्यते}\ A; \skt{अस्मिन्विंदन्त्य्}\ K\textsuperscript{pc}}}\skt{ }\edtext{\edtext{\skt{वायुर्विद्यत}}{\lemma{\skt{वायुर्}}\Afootnote{A K; \skt{वायु}\ N; \skt{चा}\ H\textsuperscript{ac}; \skt{चायु}\ H\textsuperscript{pc}}}}{\lemma{\skt{विद्यत}}\Afootnote{N; \skt{विन्दन्ति’}\ A; \skt{विन्द्यत}\ K; \emph{om.}\ H\textsuperscript{ac}; \skt{र्व्विद्यत}\ H\textsuperscript{pc}}}\skt{ }\edtext{\edtext{\skt{इत्यायुर्वेदः}}{\lemma{\skt{इत्य्}}\Afootnote{A K N H\textsuperscript{pc}; \emph{om.}\ H\textsuperscript{ac}}}}{\lemma{\skt{आयुर्वेदः}}\Afootnote{A K N H\textsuperscript{pc}; \skt{युर्व्वेदस्}\ H\textsuperscript{ac}}}\skt{ ।}\pend

\pstart \ledsidenote{\passref{1.1.16}}\edtext{\edtext{\skt{तस्याङ्गवरमागमप्रत्यक्षानुमानोपमानैरविरुद्धमुच्यमानमुपधारयध्वम्}}{\lemma{\skt{आगमप्रत्यक्षानुमानोपमानैर्}}\Afootnote{K\textsuperscript{ac} N H; \skt{आद्यं प्रत्यक्षागमानुमानोपमानैर्}\ A; \skt{आगमप्रत्यक्षानुनामानोपमानैर्}\ K\textsuperscript{pc}}}}{\lemma{\skt{उपधारयध्वम्}}\Afootnote{K N H; \skt{उपधारय}\ A}}\skt{ ।}\pend

\pstart \ledsidenote{\passref{1.1.17}}\skt{एतद्ध्यङ्गं प्रथमं }\edtext{\skt{प्रधानं}}{\Afootnote{K N H; \emph{om.}\ A}}\skt{ }\edtext{\skt{प्रागभिहितत्वाद्व्रणसंरोहणकरत्वाद्यज्ञशिरःप्रधानसन्धानाच्च}}{\Afootnote{K H; \skt{प्रागभिघातव्रणसंरोहाद्यज्ञशिरःसन्धानाच्च}\ A; \skt{प्रागभिधानत्वाद्व्रणसंरोहणकरत्वाद्यज्ञशिरप्रधानसन्धानाच्च}\ N}}\skt{ । श्रूयते हि }\edtext{\skt{यथा}}{\Afootnote{A K\textsuperscript{pc} N H; \emph{om.}\ K\textsuperscript{ac}}}\skt{ }\edtext{\skt{पुरा}}{\Afootnote{K N H; \emph{om.}\ A}}\skt{ रुद्रेण }\edtext{\edtext{\edtext{\skt{शिरश्छिन्नमश्विभ्यां}}{\lemma{\skt{शिरश्}}\Afootnote{K; \skt{यज्ञस्य शिरश्छिन्नमिति, ततो देवा अश्विनावभिगम्योचुः- भगवन्तौ नः श्रेष्ठतमौ युवां भविष्यथः, भवद्भ्यां यज्ञस्य शिरः सन्धातव्यमिति । तावूचतुरेवमस्त्विति । अथ तयोरर्थे देवा इन्द्रं यज्ञभागेन प्रासादयन्}\ A; \skt{यज्ञस्य शिरश्छिन्नमश्वि}\ N; \skt{शिरश्छिन्नम्}\ H\textsuperscript{ac}; \skt{यज्ञ शिरश्छिन्नम्}\ H\textsuperscript{pc}}}}{\lemma{\skt{छिन्नम्}}\Afootnote{K; \skt{ताभ्यां}\ A; \emph{om.}\ N H}}}{\lemma{\skt{अश्विभ्यां}}\Afootnote{K H; \skt{यज्ञस्य}\ A; \skt{भ्यां}\ N}}\skt{ }\edtext{\skt{पुनः}}{\Afootnote{K N H; \skt{शिरः}\ A}}\skt{ }\edtext{\edtext{\skt{सन्धितमिति}}{\lemma{\skt{सन्धितम्}}\Afootnote{K N H; \skt{संहितम्”}\ A}}}{\lemma{\skt{इति}}\Afootnote{A N; \skt{इत्य्}\ K H}}\skt{ ।}\pend

\pstart \ledsidenote{\passref{1.1.18}}\edtext{\edtext{\skt{अष्टानामपि}}{\lemma{\skt{अष्टानाम्}}\Afootnote{K N H; \skt{अष्टा}\ A}}}{\lemma{\skt{अपि}}\Afootnote{K N H; \skt{स्वपि}\ A}}\skt{ }\edtext{\edtext{\skt{चायुर्वेदतन्त्राणामेतदेवाधिकमाशु}}{\lemma{\skt{चायुर्वेदतन्त्राणाम्}}\Afootnote{K N H; \skt{चायुर्वेदतन्त्रे}\ A}}}{\lemma{\skt{एतदेवाधिकम्}}\Afootnote{K N H; \skt{ष्वेतदेवाधिकमभिमतम्}\ A}}\skt{ }\edtext{\skt{क्रियाकरणाद्यन्त्रशस्त्रक्षाराग्निप्रणिधानात्सर्वतन्त्रसामान्याच्च}}{\Afootnote{K N H; \skt{क्रियाकरणाद्यन्त्रशस्त्रक्षाराग्निप्रणिधानात्सर्वतन्त्रसामान्याच्च}\ A}}\skt{ ।}\pend

\pstart \ledsidenote{\passref{1.1.19}}\edtext{\skt{तदिदं}}{\Afootnote{A K H; \skt{तदिद}\ N}}\skt{ }\edtext{\skt{शाश्वतं}}{\Afootnote{A K H; \skt{शाश्वत}\ N}}\skt{ पुण्यं }\edtext{\skt{स्वर्ग्यं}}{\Afootnote{A N; \skt{स्वर्गय}\ K\textsuperscript{ac}; \skt{स्वर्ग}\ K\textsuperscript{pc}; \skt{स्वर्ग्यां यं}\ H}}\skt{ }\edtext{\skt{यशस्यमायुष्यं}}{\lemma{\skt{यशस्यम्}}\Afootnote{A K N H\textsuperscript{pc}; \skt{शस्यम्}\ H\textsuperscript{ac}}}\skt{ }\edtext{\skt{वृत्तिकरञ्च}}{\Afootnote{K N H; \skt{वृत्तिकरं चेति}\ A}}\skt{ ।}\pend

\pstart \ledsidenote{\passref{1.1.20}}\edtext{\skt{तद्ब्रह्मा}}{\Afootnote{K; \skt{ब्रह्मा}\ A; \skt{त ब्रह्मा}\ N; \skt{तद्ब्रह्मो}\ H}}\skt{ }\edtext{\skt{प्रोवाच}}{\Afootnote{A K N; \skt{वाच}\ H}}\skt{ }\edtext{\edtext{\skt{तत्प्रजापतिरधिजगे}}{\lemma{\skt{तत्प्रजापतिर्}}\Afootnote{K N H; \skt{ततः प्रजापतिर्}\ A}}}{\lemma{\skt{अधिजगे}}\Afootnote{A K H; \skt{अधिजचे}\ N}}\skt{ तस्मादश्विनावश्विभ्यामिन्द्र }\edtext{\skt{इन्द्रादहं}}{\Afootnote{A N; \skt{इन्द्रादःअं}\ K H}}\skt{ मया त्विह प्रदेयमर्थिभ्यः प्रजाहितहेतोः ।}\pend

\pstart\versepar \ledsidenote{\passref{1.1.21}}\edtext{\skt{भवति चात्र}}{\Afootnote{A; \skt{भ वति चात्र}\ K; \skt{भवति चात्र}\ N H}}\skt{ ।}\footnoteA{The Nepalese witnesses commonly use the abbreviation "\emph{bha}" for "\emph{bhavati cātra}", introducing the next verse.}\skt{ }\pend
\pstart\versepar \skt{अहं हि धन्वन्तरिरादिदेवो जरारुजामृत्युहरो }\edtext{\skt{मराणां}}{\Afootnote{K N; \skt{ऽमराणाम्}\ A H}}\skt{ । }\pend
\pstart\versepar \skt{}\edtext{\edtext{\skt{शल्यम्महच्छास्त्रवरं}}{\lemma{\skt{शल्यम्}}\Afootnote{N H; \skt{शल्याङ्}\ A; \skt{…ल्यम्}\ K}}}{\lemma{\skt{महच्छास्त्रवरं}}\Afootnote{K H; \skt{गमङ्गैरपरैर्}\ A; \skt{महच्छस्त्रवरं}\ N}}\skt{ }\edtext{\skt{गृहीत्वा}}{\Afootnote{K N H; \skt{उपेतं}\ A}}\skt{ }\edtext{\skt{प्राप्तोऽस्मि गां}}{\lemma{\skt{ऽस्मि गां}}\Afootnote{A K; \skt{मि गां}\ N; \skt{ऽस्मि गम्}\ H}}\skt{ भूय इहोपदेष्टुं ।}\pend

\pstart \ledsidenote{\passref{1.1.22}}\edtext{\skt{तत्रास्मिञ्शास्त्रे}}{\Afootnote{H; \skt{अस्मिञ्छास्त्रे}\ A; \skt{तत्रास्मिन्शास्त्रे}\ K; \skt{तत्रास्मिञ्छास्त्रे}\ N}}\skt{ पञ्चमहाभूतशरीरिसमवायः पुरुष इत्युच्यते । }\edtext{\skt{तस्मिन्क्रिया }\edtext{\skt{सोऽधिष्ठानं}}{\lemma{\skt{ऽधिष्ठानं}}\Afootnote{A K H; \skt{धिष्ठानं}\ N}}}{\lemma{\skt{तस्मिन्क्रिया सो}}\Afootnote{A K H; \skt{तस्मिं क्रिया सो}\ N}}\skt{ । }\edtext{\skt{कस्माल् । लोकद्वैविध्याल् । लोको हि}}{\Afootnote{H; \skt{कस्मात्? लोकस्य द्वैविध्यात ; लोको हि}\ A; \skt{कस्मात्लोकद्वैविध्याल्लोको हि}\ K; \skt{कस्माल्लौकिकविध्याल्लोको हि}\ N}}\skt{ }\edtext{\skt{द्विविधः}}{\Afootnote{A; \skt{द्विविध}\ K N H}}\footnoteA{The reading of \emph{dvividhan} might be a corruption of \emph{dvividhas}, which is a permitted sandhi.}\skt{ }\edtext{\skt{स्थावरो}}{\Afootnote{A; \skt{न्स्थावरो}\ K H; \skt{स्थावर}\ N}}\skt{ }\edtext{\skt{जङ्गमश्च}}{\Afootnote{K N H; \skt{जङ्गमश्च}\ A}}\skt{ । द्विविधात्मक एवाग्नेयः }\edtext{\skt{सौम्यश्च}}{\Afootnote{K N H; \skt{सौम्यश्च}\ A}}\skt{ }\edtext{\skt{तद्भूयस्त्वात्पञ्चात्मको वा}}{\Afootnote{A K H; \skt{तद्भूयस्त्वात्पञ्चातको वा}\ N}}\skt{ । तत्र चतुर्विधो भूतग्रामः }\edtext{\skt{संस्वेदजाद्रिजजरायुजाण्डजसंज्ञः}}{\Afootnote{\skt{संस्वेदजजरायुजाण्डजोद्भिज्जसञ्ज्ञः}\ A; \skt{संस्वेदजाद्रिजजरायुजाण्डजसंज्ञास्}\ K; \skt{संस्वेदजाद्रिदजजरायुजाण्डजसं…स्}\ N; \skt{संस्वेदजो ह्रिदजरायुजाण्डजसंज्ञास्}\ H}}\skt{ । }\edtext{\skt{तस्मिन्पुरुषः}}{\Afootnote{H; \skt{तत्र पुरुषः}\ A; \skt{तस्मिस्पुरुषः}\ K; \skt{तस्मिं पुरुषः}\ N}}\skt{ }\edtext{\edtext{\edtext{\skt{प्रधानस्तस्योपकरणमन्यत्}}{\lemma{\skt{प्रधानस्}}\Afootnote{K N H; \skt{प्रधानं}\ A}}}{\lemma{\skt{तस्योपकरणम्}}\Afootnote{K N H; \skt{तस्योपकरण}\ A}}\skt{ । }\edtext{\skt{तस्मात्पुरुषोऽधिष्ठानं}}{\lemma{\skt{ऽधिष्ठानं}}\Afootnote{A H; \skt{ऽधिष्ठान}\ K; \skt{धिष्ठानं}\ N}}}{\lemma{\skt{अन्यत् । तस्मात्पुरुषो}}\Afootnote{K H; \skt{मन्यत्; तस्मात्पुरुषो}\ A; \skt{अन्यं तस्मात्पुरुषो}\ N}}\skt{ ।}\pend

\pstart \ledsidenote{\passref{1.1.23}}\edtext{\skt{तद्दुःखसंयोगा}}{\Afootnote{A; \skt{तद्दुःखसंयोगा}\ K N H}}\skt{ व्याधय }\edtext{\skt{इत्युच्यन्ते}}{\lemma{\skt{इत्य्}}\Afootnote{K N H; \emph{om.}\ A}}\skt{ ।}\pend

\pstart \ledsidenote{\passref{1.1.24}}\skt{ते चतुर्विधा आगन्तवः शारीरा }\edtext{\skt{मानसा}}{\Afootnote{K N H; \skt{मानसाः}\ A}}\skt{ }\edtext{\edtext{\skt{स्वाभाविकाश्चेति}}{\lemma{\skt{स्वाभाविकाश्}}\Afootnote{K N H; \skt{स्वाभाविका}\ A}}}{\lemma{\skt{चेति}}\Afootnote{K N H; \skt{श्चेति}\ A}}\skt{ ।}\pend

\pstart \ledsidenote{\passref{1.1.25.1}}\edtext{\edtext{\skt{तेष्वागन्तवोऽभिघातनिमित्ताः}}{\lemma{\skt{तेष्व्}}\Afootnote{K N H; \skt{तेषां}\ A}}}{\lemma{\skt{ऽभिघातनिमित्ताः}}\Afootnote{A H; \skt{ऽभिघातनिमित्ता}\ K; \skt{भिघातनिमित्ता}\ N}}\skt{ ॥}\pend

\pstart \ledsidenote{\passref{1.1.25.2}}\edtext{\skt{शारीरास्त्वन्नमूला}}{\lemma{\skt{त्वन्नमूला}}\Afootnote{K N H; \skt{त्वन्नपानमूला}\ A}}\skt{ }\edtext{\skt{वातपित्तकफशोणितवैषम्यनिमित्ताः}}{\Afootnote{K N H; \skt{वातपित्तकफशोणितसन्निपातवैषम्यनिमित्ताः}\ A}}\skt{ ।}\pend

\pstart \ledsidenote{\passref{1.1.25.3}}\skt{मानसास्तु }\edtext{\skt{क्रोधशोकदैन्यहर्षकामविषादेर्ष्यासूयामात्सर्यलोभादय}}{\Afootnote{N H; \skt{क्रोधशोकभयहर्षविषादेर्ष्याभ्यसूयादैन्यमात्सर्यकामलोभप्रभृतय}\ A; \skt{क्रोधाशोकदैन्यहर्षकामविषादेर्ष्यासूयामात्सर्यलोभादय}\ K}}\skt{ }\edtext{\skt{इच्छाद्वेषनिमित्ताः}}{\Afootnote{K N H; \skt{इच्छाद्वेषभेदैर्भवन्ति}\ A}}\skt{ ।}\pend

\pstart \ledsidenote{\passref{1.1.25.4}}\skt{स्वाभाविकास्तु }\edtext{\skt{क्षुत्पिपासाजरामृत्युनिद्राप्रकृतयः}}{\Afootnote{K H; \skt{क्षुत्पिपासाजरामृत्युनिद्राप्रकृ(भृ)तयः}\ A; \skt{क्षुत्पिपासाजरामृत्युनिद्राप्रभृतयः}\ N}}\skt{ ।}\pend

\pstart \ledsidenote{\passref{1.1.26}}\edtext{\skt{त एते}}{\Afootnote{A K H; \skt{एते}\ N}}\skt{ }\edtext{\skt{मनःशरीराधिष्ठाना}}{\Afootnote{N H; \skt{मनःशरीराधिष्ठानाः}\ A; \skt{मनःशरीराधि…}\ K}}\skt{ }\edtext{\skt{भवन्ति}}{\Afootnote{K N H; \emph{om.}\ A}}\skt{ ।}\pend

\pstart \ledsidenote{\passref{1.1.27}}\skt{तेषां }\edtext{\skt{लेखनबृम्हणसंशोधनसंशमनाहाराचाराः}}{\Afootnote{K N; \skt{संशोधनसंशमनाहाराचाराः}\ A; \skt{लेखनबृम्हनसंशोधनसंशमनाहाराचाराः}\ H}}\skt{ सम्यक्प्रयुक्ता निग्रहहेतवो }\edtext{\skt{भवन्ति}}{\Afootnote{K N H; \emph{om.}\ A}}\skt{ ।}\pend

\pstart \ledsidenote{\passref{1.1.28}}\skt{प्राणिनां }\edtext{\skt{पुनर्मूलमाहारो}}{\lemma{\skt{पुनर्}}\Afootnote{A K H; \skt{पुन}\ N}}\skt{ बलवर्णौजसाञ्च । स षट्सु }\edtext{\skt{रसेष्वायत्तः । रसाः}}{\lemma{\skt{आयत्तः । रसाः}}\Afootnote{A; \skt{आयत्तरसाः}\ K N H}}\skt{ }\edtext{\skt{पुनर्द्रव्याश्रयिनः}}{\lemma{\skt{द्रव्याश्रयिनः}}\Afootnote{K; \skt{द्रव्याश्रयाः}\ A; \skt{दव्याश्रयिना}\ N; \skt{द्द्रव्याश्रयिणः}\ H}}\skt{ । द्रव्याणि }\edtext{\skt{पुनरोषध्यः । ता}}{\lemma{\skt{ओषध्यः । ता}}\Afootnote{K N; \skt{ओषधयः । तास्तु}\ A; \skt{औषध्यः । ता}\ H}}\skt{ }\edtext{\skt{द्विविधा}}{\Afootnote{K N; \skt{द्विविधाः}\ A H}}\skt{ । स्थावरा जङ्गमाश्च ।}\pend

\pstart \ledsidenote{\passref{1.1.29}}\skt{तासां स्थावराश्चतुर्विधा वनस्पतयो वृक्षा }\edtext{\skt{ओषध्यो}}{\Afootnote{K H; \emph{om.}\ A; \skt{औषध्यो}\ N}}\skt{ }\edtext{\skt{वीरुध}}{\Afootnote{K N H; \skt{वीरुध, ओषधय}\ A}}\skt{ इति । }\edtext{\skt{तास्वपुष्पा}}{\lemma{\skt{अपुष्पा}}\Afootnote{K N H; \skt{अपुष्पाः}\ A}}\skt{ }\edtext{\skt{फलवन्त्यो}}{\Afootnote{K H; \skt{फलवन्तो}\ A; \skt{…लवन्त्यो}\ N}}\skt{ वनस्पतयः । }\edtext{\skt{पुष्पफलो ये ता}}{\Afootnote{K H; \skt{पुष्पफलवन्तो}\ A; \skt{पुष्पफलो य्ये ता}\ N}}\skt{ वृक्षाः । }\edtext{\edtext{\skt{फलपाकनिष्ठास्त्वोषध्यः}}{\lemma{\skt{फलपाकनिष्ठास्}}\Afootnote{K N H; \skt{प्रतानवत्यः स्तम्बिन्यश्च वीरुधः, फलपाकनिष्ठा}\ A}}}{\lemma{\skt{त्वोषध्यः}}\Afootnote{K H; \skt{ओष}\ A; \skt{त्वौषध्यः}\ N}}\skt{ । }\edtext{\skt{प्रतानवत्यो}}{\Afootnote{H; \skt{ध}\ A; \skt{प्रतानव…}\ K; \skt{प्रतानवन्त्यो}\ N}}\skt{ }\edtext{\skt{वीरुध}}{\Afootnote{K H; \skt{य}\ A; \skt{विरुध}\ N}}\skt{ इति ।}\pend

\pstart \ledsidenote{\passref{1.1.30}}\edtext{\skt{जङ्गमाः}}{\Afootnote{K H; \skt{जङ्गमा}\ A N}}\skt{ खल्वपि }\edtext{\skt{चतुर्विधाः}}{\Afootnote{A N; \skt{चतुविधा}\ K H}}\skt{ । }\edtext{\skt{जरायुजाण्डजसंस्वेदजोद्भिदा}}{\Afootnote{H; \skt{जरायुजाण्डजस्वेदजोद्भि}\ A; \skt{जरायुजा…जस…स्वेदजोद्भिदाः}\ K; \skt{जरायुजाण्डजसस्वेदजोद्भिदाः}\ N}}\skt{ }\edtext{\skt{इति}}{\Afootnote{K N H; \skt{ज्जाः}\ A}}\skt{ । }\edtext{\skt{तेषाम्पशुमनुष्यव्यालादयो}}{\lemma{\skt{तेषाम्}}\Afootnote{K N H; \skt{तत्र}\ A}}\skt{ जरायुजाः । }\edtext{\edtext{\skt{खगसरीसृपसर्पास्त्वण्डजाः}}{\lemma{\skt{खगसरीसृपसर्पास्}}\Afootnote{N H; \skt{खगसर्पसरीसृपप्रभृतयो}\ A; \skt{…सृप स्}\ K\textsuperscript{ac}; \skt{…सृप सर्पा स्}\ K\textsuperscript{pc}}}}{\lemma{\skt{त्वण्डजाः}}\Afootnote{K N H; \skt{ऽण्डजाः}\ A}}\skt{ । }\edtext{\skt{कृमिकुन्थपिपीलिकाप्रभृतयः}}{\Afootnote{K; \skt{कृमिकीटपिपीलिकाप्रभृतयः}\ A; \skt{कृमिकुट्ठपिपीलिकाप्रभृतयः}\ N; \skt{क्रिमिकुष्ठपिपीलिकाप्रभृतयः}\ H}}\footnoteA{\emph{kuṭṭha} is Prākṛta for "skin disease, leprosy," but this sense does not fit this passage; the word here must refer to an insect or similar creature. In our MSS, \emph{kuṭṭha} appears to be a transmitted form, possibly Sanskritized to \emph{kuṣṭha} by the scribe of MS H. It is not obvious how it could be a scribal error for the Vulgate's "\emph{kīṭa}." We read with K.}\skt{ }\edtext{\skt{संस्वेदजाः}}{\Afootnote{K N H; \skt{स्वेदजाः}\ A}}\skt{ । }\edtext{\edtext{\skt{इन्द्रगोपकमण्डूकप्रभृतयश्चोद्भिदाः}}{\lemma{\skt{इन्द्रगोपकमण्डूकप्रभृतयश्}}\Afootnote{K N H; \skt{इन्द्रगोपमण्डूकप्रभृतय}\ A}}}{\lemma{\skt{चोद्भिदाः}}\Afootnote{K N H; \skt{उद्भिज्जाः}\ A}}\skt{ ।}\pend

\pstart \ledsidenote{\passref{1.1.31}}\skt{तत्र }\edtext{\skt{स्थावरेभ्यस्त्वक्पत्रपुष्पफलमूलकन्दक्षीरनिर्याससारस्नेहस्वरसाः}}{\lemma{\skt{त्वक्पत्रपुष्पफलमूलकन्दक्षीरनिर्याससारस्नेहस्वरसाः}}\Afootnote{K H; \skt{त्वक्पत्रपुष्पफलमूलकन्दनिर्यासस्वरसादयः}\ A; \skt{त्वक्पत्रपुष्पफलमूलकन्दक्षीरनिर्याससारस्नेहस्वरसा}\ N}}\skt{ प्रयोजनवन्तः । }\edtext{\skt{जङ्गमेभ्यश्चर्मरोमनखरुधिरादयः}}{\lemma{\skt{चर्मरोमनखरुधिरादयः}}\Afootnote{K N H; \skt{चर्मनखरोमरुधिरादयः}\ A}}\skt{ ।}\pend

\pstart \ledsidenote{\passref{1.1.32}}\edtext{\skt{पार्थिवस्तु}}{\Afootnote{K N H; \skt{पार्थिवाः}\ A}}\skt{ }\edtext{\skt{सुवर्णरजतादयः}}{\Afootnote{K N H; \skt{सुवर्णरजतमणिमुक्तामनःशिलामृत्कपालादयः}\ A}}\skt{ ।}\pend

\pstart \ledsidenote{\passref{1.1.33}}\edtext{\skt{कालकृतास्तु}}{\Afootnote{K N H; \skt{कालकृताः}\ A}}\skt{ }\edtext{\skt{प्रवातनिवातातपच्छायातमोज्योत्स्नाशीतोष्णवर्षासु}}{\Afootnote{K N H; \skt{प्रवातनिवातातपच्छायाज्योत्स्नातमःशीतोष्णवर्षा}\ A}}\skt{ }\edtext{\skt{सम्प्लवाः}}{\Afootnote{K H; \emph{om.}\ A; \skt{ससम्प्लवाः}\ N}}\skt{ । }\edtext{\skt{कालविशेषास्तु}}{\Afootnote{K N H; \emph{om.}\ A}}\skt{ }\edtext{\skt{निमेषकाष्ठाकलामुहूर्ताहोरात्रपक्षमासर्त्वयनसम्वत्सरयुगविशेषाः}}{\Afootnote{K N H; \skt{होरात्रपक्षमासर्त्वयनादयः संवत्सरविशेषाः}\ A}}\skt{ ।}\pend

\pstart \ledsidenote{\passref{1.1.34}}\edtext{\skt{स्वभावत}}{\Afootnote{K N H; \skt{त एते स्वभावत}\ A}}\skt{ एव दोषाणां }\edtext{\skt{सञ्चयप्रकोपोपशमप्रतीकारहेतवो}}{\Afootnote{K H; \skt{सञ्चयप्रकोपप्रशमप्रतीकारहेतवः}\ A; \skt{सञ्चयप्रकोपशमप्रतीकारहेतवो}\ N}}\skt{ }\edtext{\skt{भवन्ति}}{\Afootnote{K N H; \emph{om.}\ A}}\skt{ । प्रयोजनवन्तश्च ।}\pend

\pstart\versepar \ledsidenote{\passref{1.1.35}}\skt{भवन्ति चात्र । }\pend
\pstart\versepar \skt{शारीराणां विकाराणामेष }\edtext{\edtext{\skt{वर्गश्चतुर्विधः}}{\lemma{\skt{वर्गश्}}\Afootnote{A K H; \skt{वर्गेश्}\ N}}}{\lemma{\skt{चतुर्विधः}}\Afootnote{A N H; \skt{चतुविधः}\ K}}\skt{ । }\pend
\pstart\versepar \skt{प्रकोपे प्रशमे चैव }\edtext{\edtext{\skt{हेतुरुक्तश्चिकित्सकैः}}{\lemma{\skt{उक्तश्}}\Afootnote{K H; \skt{उक्त}\ A N}}}{\lemma{\skt{चिकित्सकैः}}\Afootnote{K N H; \skt{श्चिकित्सकैः}\ A}}\skt{ ॥}\edtext{}{\lemma{\skt{भवन्ति चात्र । शारीराणां}}\Afootnote{\skt{भवन्ति चात्र श्लोकाः- शारीराणां}\ A; \skt{भवन्ति चात्र॥ शारीराणां}\ K; \skt{भ शारीराणां}\ N H}}\pend

\pstart\versepar \ledsidenote{\passref{1.1.36}}\skt{आगन्तवस्तु ये रोगास्ते द्विधा }\edtext{\skt{निपतन्ति ह}}{\Afootnote{K N H; \skt{निपतन्ति हि}\ A}}\skt{ । }\pend
\pstart\versepar \skt{मनस्यन्ये }\edtext{\skt{शरीरे च}}{\Afootnote{K N H; \skt{शरीरेऽन्ये}\ A}}\skt{ तेषान्तु द्विविधा क्रिया ॥}\pend

\pstart\versepar \ledsidenote{\passref{1.1.37}}\skt{शरीरपतितानां तु }\edtext{\skt{शारीरवदुपक्रमः}}{\Afootnote{K N H; \skt{शारीरवदुपक्रमः}\ A}}\skt{ । }\pend
\pstart\versepar \skt{मानसानान्तु }\edtext{\edtext{\skt{शब्दादिरिष्टो}}{\lemma{\skt{शब्दादिर्}}\Afootnote{K N H; \skt{शब्दादि}\ A}}}{\lemma{\skt{इष्टो}}\Afootnote{K N H; \skt{रिष्टो}\ A}}\skt{ वर्गः सुखावहः ॥}\footnoteA{Note the good example of the inconsistent use by the scribe of witness H of anusvāra in 37ab (\emph{°n tu}) vs 37cd (\emph{°ṃ tu}).}\pend

\pstart \ledsidenote{\passref{1.1.38}}\skt{एवमेतत्पुरुषो व्याधिरौषधं क्रिया काल इति }\edtext{\skt{समासेन}}{\Afootnote{K N H; \emph{om.}\ A}}\skt{ चतुष्टयं }\edtext{\edtext{\skt{व्याख्यातम्भवति}}{\lemma{\skt{व्याख्यातम्}}\Afootnote{K N H; \skt{समासेन व्याख्यातम्}\ A}}}{\lemma{\skt{भवति}}\Afootnote{K N H; \emph{om.}\ A}}\skt{ । तत्र पुरुषग्रहणात्तत्संभवद्रव्यसमूहो भूतादिरुक्तः }\edtext{\skt{तदङ्गप्रत्यङ्गविकल्पाश्च}}{\Afootnote{A K; \skt{तदगद्गप्रत्यङ्गविकल्पाश्च}\ N; \skt{तदङ्गाप्रत्यङ्गा विकल्पाश्च}\ H\textsuperscript{ac}; \skt{तदङ्गप्रत्यङ्गा विकल्पाश्च}\ H\textsuperscript{pc}}}\skt{ }\edtext{\skt{त्वङ्मांसासिरास्नाय्वस्थिसन्धिप्रभृतयः}}{\Afootnote{K N H; \skt{त्वङ्मांसास्थिसिरास्नायुप्रभृतयः}\ A}}\skt{ । }\edtext{\skt{व्याधिग्रहणाद्वातपित्तकफशोणितसन्निपातागन्तुस्वभावनिमित्ताः}}{\Afootnote{K H; \skt{व्याधिग्रहणाद्वातपित्तकफशोणितसन्निपातवैषम्यनिमित्ताः}\ A; \skt{व्याधिग्राहणाद्वातपित्तकफशोणितसन्निपातागन्तुस्वभावनिमित्ताः}\ N}}\skt{ सर्व एव व्याधयो }\edtext{\skt{व्याख्याता}}{\Afootnote{K N H; \skt{व्याख्याताः}\ A}}\skt{ }\edtext{\skt{भवन्ति}}{\Afootnote{K H; \emph{om.}\ A; \skt{भवस्वन्ति}\ N}}\skt{ । }\edtext{\edtext{\skt{ओषधग्रहणाद्द्रव्यरसगुणवीर्यविपाकानामादेशः}}{\lemma{\skt{ओषधग्रहणाद्द्रव्यरसगुणवीर्यविपाकानाम्}}\Afootnote{K N; \skt{औषधग्रहणाद्द्रव्यरसगुणवीर्यविपाकाना}\ A; \skt{ओषधग्रहणाद्द्रव्यरसगुणवीर्यविपाको ना}\ H}}}{\lemma{\skt{आदेशः}}\Afootnote{K N; \skt{मादेशः}\ A H}}\skt{ । }\edtext{\skt{क्रियाग्रहणात्स्नेहादीनि}}{\Afootnote{A K N H\textsuperscript{pc}; \skt{क्रियाग्रहणात्स्नेहा}\ H\textsuperscript{ac}}}\skt{ }\edtext{\skt{च्छेद्यादीनि च}}{\Afootnote{A K N H\textsuperscript{pc}; \skt{दीनि च}\ H\textsuperscript{ac}}}\skt{ }\edtext{\skt{कर्माण्युपदिष्टानि}}{\lemma{\skt{उपदिष्टानि}}\Afootnote{K N H; \skt{व्या}\ A}}\skt{ }\edtext{\skt{भवन्ति}}{\Afootnote{K N H; \skt{ख्यातानि}\ A}}\skt{ । कालग्रहणात्सर्व }\edtext{\skt{एव}}{\Afootnote{K N H; \emph{om.}\ A}}\skt{ }\edtext{\skt{क्रियाकालादेशः}}{\Afootnote{K; \skt{क्रियाकालानामादेशः}\ A; \skt{क्रियाकालनिर्देशः}\ N; \skt{क्रियाकालकालदेशः}\ H}}\skt{ ।}\pend

\pstart\versepar \ledsidenote{\passref{1.1.39}}\edtext{\skt{भवति चात्र}}{\Afootnote{A; \emph{om.}\ K; \skt{भवति चात्र}\ N H}}\skt{ }\pend
\pstart\versepar \skt{बीजञ्चिकित्सितस्यैतत्समासेन }\edtext{\skt{प्रकीतम्}}{\Afootnote{\skt{प्रकीर्तितम्}\ A K N H}}\skt{ । }\pend
\pstart\versepar \skt{}\edtext{\edtext{\edtext{\skt{सविंशमध्यायशतमस्य}}{\lemma{\skt{सविंशम्}}\Afootnote{K N H; \skt{सविंश}\ A}}}{\lemma{\skt{अध्यायशतम्}}\Afootnote{K N H; \skt{मध्यायशत}\ A}}}{\lemma{\skt{अस्य}}\Afootnote{K N H; \skt{मस्य}\ A}}\skt{ }\edtext{\skt{व्याख्या}}{\Afootnote{A K\textsuperscript{pc} N H; \emph{om.}\ K\textsuperscript{ac}}}\skt{ भविष्यति ॥}\pend

\pstart \ledsidenote{\passref{1.1.40}}\edtext{\skt{तच्च}}{\Afootnote{A H; \skt{त}\ K N}}\skt{ }\edtext{\skt{विंशमध्यायशतं}}{\lemma{\skt{विंशम्}}\Afootnote{H; \skt{सविंशम्}\ A; \skt{च्च विंशम्}\ K N}}\skt{ पञ्चसु स्थानेषु }\edtext{\skt{चेति}}{\Afootnote{K N H; \emph{om.}\ A}}\skt{ । }\edtext{\skt{तत्र}}{\Afootnote{A K N; \skt{अत्र}\ H}}\skt{ }\edtext{\edtext{\skt{श्लोकस्थाननिदानशारीरचिकित्सितकल्पेष्वर्थवशाद्विभज्य}}{\lemma{\skt{श्लोकस्थाननिदानशारीरचिकित्सितकल्पेष्व्}}\Afootnote{N H; \skt{सूत्रनिदानशारीरचिकित्सितकल्पेष्व्}\ A; \skt{श्लोकस्थाननिदानशारी}\ K}}\skt{ उत्तरे}}{\lemma{\skt{अर्थवशाद्विभज्य उत्तरे}}\Afootnote{\skt{अर्थवशात्संविभज्योत्तरे}\ A; \skt{…}\ K; \skt{अर्थवशा}\ N; \skt{अर्थवसाद्विभज्य उत्तरे}\ H}}\skt{ }\edtext{\skt{वक्ष्यामः}}{\Afootnote{H; \skt{तन्त्रे शेषानर्थान्वक्ष्यामः}\ A; \skt{…इभज्य उत्तरे वक्ष्यामः}\ K; \skt{द्विभज्योत्तरे वक्ष्यामः}\ N}}\skt{ ।}\pend

\pstart\versepar \ledsidenote{\passref{1.1.41}}\edtext{\skt{भवति चात्र}}{\Afootnote{A; \emph{om.}\ K; \skt{भवति चात्र}\ N H}}\skt{ । }\pend
\pstart\versepar \skt{स्वयम्भुवा प्रोक्तमिदं सनातनम् }\pend
\pstart\versepar \skt{}\edtext{\skt{पठेत्तु यः}}{\Afootnote{K N H; \skt{पठेद्धि यः}\ A}}\skt{ काशिपतिप्रकाशितम् । }\pend
\pstart\versepar \skt{स पुण्यकर्मा भुवि पूजितो नृपैर् }\pend
\pstart\versepar \skt{असुक्षये }\edtext{\skt{शक्रसलोकतामियादिति}}{\lemma{\skt{इयादिति}}\Afootnote{\skt{व्रजेत्}\ A; \skt{इयाद्}\ K N H}}\skt{ ॥}\pend

\pstart \ledsidenote{\passref{1.1.trailer}}\edtext{\skt{॥ १ ॥}}{\Afootnote{\skt{इति सुश्रुतसंहितायां सूत्रस्थाने वेदोत्पत्तिर्नाम प्रथमोऽध्यायः}\ A; \skt{इति ॥{\symfont ❈}}\ K; \skt{इति}\ N H}}\pend

\endnumbering
